\begingroup
\let\cleardoublepage\relax


\chapter*{Java Datentypen}
\label{ch:java-datentypen}
\addcontentsline{toc}{chapter}{Java Datentypen}
\endgroup


\section*{Primitive Datentypen}
\addcontentsline{toc}{section}{Primitive Datentypen}

Java unterscheidet zwischen primitiven Datentypen und Referenztypen (Objekten).
Primitive Datentypen sind die grundlegenden Bausteine und speichern direkt Werte.

% ============================================================
\newpage


\subsection*{int - Ganzzahlen}
\addcontentsline{toc}{subsection}{int - Ganzzahlen}

\subsubsection*{Was ist int?}\\
Der Datentyp \texttt{int} (Integer) speichert ganze Zahlen ohne Nachkommastellen.

\subsubsection*{Wertebereich:}
\begin{itemize}
	\item Von: -2.147.483.648
	\item Bis: 2.147.483.647
	\item Speichergröße: 32 Bit (4 Bytes)
\end{itemize}

\subsubsection*{Syntax:}
\begin{lstlisting}
int variablenName = wert;
\end{lstlisting}

\subsubsection*{Beispiele:}
\begin{lstlisting}
// Deklaration und Initialisierung
int alter = 25;
int temperatur = -5;
int punkte = 0;

// Ohne Initialisierung
int zahl;
zahl = 42;

// Berechnungen
int summe = 10 + 20;          // summe = 30
int differenz = 50 - 15;      // differenz = 35
int produkt = 6 * 7;          // produkt = 42
int quotient = 20 / 4;        // quotient = 5
int rest = 17 % 5;            // rest = 2

// Mit Variablen rechnen
int a = 10;
int b = 3;
int ergebnis = a + b;         // ergebnis = 13

// Inkrement und Dekrement
int zaehler = 5;
zaehler++;                    // zaehler = 6
zaehler--;                    // zaehler = 5

// Zusammengesetzte Operatoren
int wert = 10;
wert += 5;                    // wert = 15 (entspricht: wert = wert + 5)
wert -= 3;                    // wert = 12
wert *= 2;                    // wert = 24
wert /= 4;                    // wert = 6
\end{lstlisting}

\begin{achtung}
	\begin{itemize}
		\item Division von int durch int ergibt int: \texttt{5 / 2 = 2} (nicht 2.5!)
		\item Bei Überschreitung des Wertebereichs: \textbf{Overflow} (Wert springt um)
		\item Standard-Wert bei Arrays: \texttt{0}
		\item Keine Nachkommastellen möglich
	\end{itemize}
\end{achtung}

% ============================================================
\newpage


\subsection*{double - Fließkommazahlen}
\addcontentsline{toc}{subsection}{double - Fließkommazahlen}

\subsubsection*{Was ist double?}\\
Der Datentyp \texttt{double} speichert Zahlen mit Nachkommastellen (Gleitkommazahlen mit doppelter Genauigkeit).

\subsubsection*{Wertebereich:}
\begin{itemize}
	\item Sehr große und sehr kleine Zahlen mit ca. 15 Dezimalstellen Genauigkeit
	\item Von: $\pm$4.9 $\times$ 10$^{-324}$
	\item Bis: $\pm$1.7 $\times$ 10$^{308}$
	\item Speichergröße: 64 Bit (8 Bytes)
\end{itemize}

\subsubsection*{Syntax:}
\begin{lstlisting}
double variablenName = wert;
\end{lstlisting}

\subsubsection*{Beispiele:}
\begin{lstlisting}
// Deklaration und Initialisierung
double pi = 3.14159;
double temperatur = -5.5;
double gewicht = 75.8;
double ergebnis = 0.0;

// Berechnungen
double summe = 10.5 + 20.3;        // summe = 30.8
double produkt = 2.5 * 4.0;        // produkt = 10.0
double quotient = 7.0 / 2.0;       // quotient = 3.5

// Wichtig: Division mit Nachkommastellen
double a = 5.0;
double b = 2.0;
double teil = a / b;               // teil = 2.5

// int zu double konvertieren
int ganzzahl = 5;
double kommazahl = ganzzahl / 2.0; // kommazahl = 2.5
// ODER
double kommazahl2 = (double) ganzzahl / 2; // kommazahl2 = 2.5

// Wissenschaftliche Notation
double klein = 1.5e-4;             // = 0.00015
double gross = 3.2e8;              // = 320000000.0

// Mathematische Funktionen
double wurzel = Math.sqrt(16.0);   // wurzel = 4.0
double potenz = Math.pow(2, 3);    // potenz = 8.0
double absolut = Math.abs(-5.5);   // absolut = 5.5
\end{lstlisting}

\begin{achtung}
	\begin{itemize}
		\item Punkt als Dezimaltrennzeichen, \textbf{nicht} Komma: \texttt{3.14} nicht \texttt{3,14}
		\item Rundungsfehler möglich: \texttt{0.1 + 0.2 = 0.30000000000000004}
		\item Division \texttt{5 / 2} ergibt \texttt{2}, aber \texttt{5.0 / 2.0} ergibt \texttt{2.5}
		\item Standard-Wert bei Arrays: \texttt{0.0}
		\item Für Geldbeträge besser \texttt{BigDecimal} verwenden
	\end{itemize}
\end{achtung}

% ============================================================
\newpage


\subsection*{boolean - Wahrheitswerte}
\addcontentsline{toc}{subsection}{boolean - Wahrheitswerte}

\subsubsection*{Was ist boolean?}\\
Der Datentyp \texttt{boolean} speichert Wahrheitswerte: entweder \texttt{true} (wahr) oder \texttt{false} (falsch).

\subsubsection*{Wertebereich:}
\begin{itemize}
	\item Nur zwei mögliche Werte: \texttt{true} oder \texttt{false}
	\item Speichergröße: 1 Bit (theoretisch), praktisch abhängig von JVM
\end{itemize}

\subsubsection*{Syntax:}
\begin{lstlisting}
boolean variablenName = true/false;
\end{lstlisting}

\subsubsection*{Beispiele:}
\begin{lstlisting}
// Deklaration und Initialisierung
boolean istWahr = true;
boolean istFalsch = false;
boolean hatBestanden = true;

// Vergleichsoperatoren ergeben boolean
int alter = 18;
boolean istErwachsen = (alter >= 18);      // true
boolean istKind = (alter < 12);            // false
boolean istGenau18 = (alter == 18);        // true
boolean istNicht18 = (alter != 18);        // false

// Logische Operatoren
boolean a = true;
boolean b = false;

boolean und = a && b;                      // false (beide muessen true sein)
boolean oder = a || b;                     // true (mindestens einer true)
boolean nicht = !a;                        // false (Negation)

// In Bedingungen
int punkte = 85;
boolean bestanden = punkte >= 50;
if (bestanden) {
    System.out.println("Bestanden!");
}

// Kombinierte Bedingungen
int temperatur = 25;
boolean angenehm = temperatur >= 18 && temperatur <= 28;

// Boolean in Schleifen
boolean weitermachen = true;
while (weitermachen) {
    // Code
    weitermachen = false;  // Schleife beenden
}

// Vergleich mit Strings
String name = "Max";
boolean istMax = name.equals("Max");       // true
\end{lstlisting}

\begin{achtung}
	\begin{itemize}
		\item Nur \texttt{true} oder \texttt{false}, \textbf{keine} Zahlen (nicht wie in C: 0 = false, 1 = true)
		\item Kleinschreibung: \texttt{true}/\texttt{false}, nicht \texttt{True}/\texttt{False}
		\item Standard-Wert bei Arrays: \texttt{false}
		\item \texttt{==} für primitive Typen, \texttt{.equals()} für Objekte (z.B. Strings)
	\end{itemize}
\end{achtung}

% ============================================================
\newpage


\subsection*{char - Einzelnes Zeichen}
\addcontentsline{toc}{subsection}{char - Einzelnes Zeichen}

\subsubsection*{Was ist char?}\\
Der Datentyp \texttt{char} speichert ein einzelnes Zeichen (Character) im Unicode-Format.

\subsubsection*{Wertebereich:}
\begin{itemize}
	\item Alle Unicode-Zeichen (Buchstaben, Ziffern, Sonderzeichen)
	\item Von: 0 bis 65.535 (als Zahlenwert)
	\item Speichergröße: 16 Bit (2 Bytes)
\end{itemize}

\subsubsection*{Syntax:}
\begin{lstlisting}
char variablenName = 'Zeichen';  // EINFACHE Anfuehrungszeichen!
\end{lstlisting}

\subsubsection*{Beispiele:}
\begin{lstlisting}
// Deklaration und Initialisierung
char buchstabe = 'A';
char ziffer = '5';
char sonderzeichen = '!';
char leerzeichen = ' ';

// Mit Unicode-Werten
char euro = '\u20AC';                      // €-Symbol
char herzchen = '\u2764';                  // ❤

// Escape-Sequenzen
char neuzeile = '\n';                      // Zeilenumbruch
char tab = '\t';                           // Tabulator
char backslash = '\\';                     // \
char anfuehrung = '\'';                    // '

// Zeichen vergleichen
char note = 'A';
if (note == 'A') {
    System.out.println("Sehr gut!");
}

// Zeichen in Berechnungen (ASCII/Unicode-Werte)
char a = 'A';                              // Unicode-Wert: 65
char b = 'B';                              // Unicode-Wert: 66
int differenz = b - a;                     // differenz = 1

// Kleinbuchstabe zu Großbuchstabe
char klein = 'a';                          // Unicode-Wert: 97
char gross = (char) (klein - 32);          // gross = 'A' (Wert: 65)

// Prüfen ob Buchstabe
char zeichen = 'M';
boolean istBuchstabe = (zeichen >= 'A' && zeichen <= 'Z') ||
                       (zeichen >= 'a' && zeichen <= 'z');

// Prüfen ob Ziffer
char z = '7';
boolean istZiffer = (z >= '0' && z <= '9');

// char aus String extrahieren
String text = "Hallo";
char erstes = text.charAt(0);              // erstes = 'H'
\end{lstlisting}

\begin{achtung}
	\begin{itemize}
		\item \textbf{Einfache} Anführungszeichen: \texttt{'A'}, nicht \texttt{"A"}
		\item Nur \textbf{ein} Zeichen: \texttt{'A'} ist char, \texttt{"A"} ist String
		\item \texttt{char} ist \textbf{nicht} dasselbe wie \texttt{String}
		\item Standard-Wert bei Arrays: \texttt{'\textbackslash u0000'} (Null-Zeichen)
		\item Kann als Zahlenwert verwendet werden (Unicode-Wert)
	\end{itemize}
\end{achtung}

% ============================================================
\newpage


\subsection*{byte - Sehr kleine Ganzzahlen}
\addcontentsline{toc}{subsection}{byte - Sehr kleine Ganzzahlen}

\subsubsection*{Was ist byte?}\\
Der Datentyp \texttt{byte} speichert sehr kleine ganze Zahlen.
Wird verwendet, wenn Speicherplatz wichtig ist.

\subsubsection*{Wertebereich:}
\begin{itemize}
	\item Von: -128
	\item Bis: 127
	\item Speichergröße: 8 Bit (1 Byte)
\end{itemize}

\subsubsection*{Syntax:}
\begin{lstlisting}
byte variablenName = wert;
\end{lstlisting}

\subsubsection*{Beispiele:}
\begin{lstlisting}
// Deklaration und Initialisierung
byte alter = 25;
byte temperatur = -15;
byte prozent = 100;

// Häufig bei Datei-Operationen
byte[] daten = new byte[1024];

// Konvertierung von int zu byte
int gross = 100;
byte klein = (byte) gross;

// Vorsicht bei Überschreitung
byte max = 127;
byte ueberlauf = (byte) (max + 1);         // ueberlauf = -128 (Overflow!)
\end{lstlisting}

\begin{achtung}
	\begin{itemize}
		\item Sehr kleiner Wertebereich: nur -128 bis 127
		\item Bei Überschreitung: Overflow
		\item Wird oft für Binärdaten verwendet
		\item Standard-Wert bei Arrays: \texttt{0}
	\end{itemize}
\end{achtung}

% ============================================================
\newpage


\subsection*{short - Kleine Ganzzahlen}
\addcontentsline{toc}{subsection}{short - Kleine Ganzzahlen}

\subsubsection*{Was ist short?}\\
Der Datentyp \texttt{short} speichert kleine ganze Zahlen.
Ein Mittelweg zwischen \texttt{byte} und \texttt{int}.

\subsubsection*{Wertebereich:}
\begin{itemize}
	\item Von: -32.768
	\item Bis: 32.767
	\item Speichergröße: 16 Bit (2 Bytes)
\end{itemize}

\subsubsection*{Syntax:}
\begin{lstlisting}
short variablenName = wert;
\end{lstlisting}

\subsubsection*{Beispiele:}
\begin{lstlisting}
// Deklaration und Initialisierung
short jahr = 2024;
short distanz = 5000;

// Konvertierung
int gross = 1000;
short klein = (short) gross;
\end{lstlisting}

\begin{achtung}
	\begin{itemize}
		\item Wird selten verwendet, meist wird direkt \texttt{int} genommen
		\item Wertebereich: -32.768 bis 32.767
		\item Standard-Wert bei Arrays: \texttt{0}
	\end{itemize}
\end{achtung}

% ============================================================
\newpage


\subsection*{long - Sehr große Ganzzahlen}
\addcontentsline{toc}{subsection}{long - Sehr große Ganzzahlen}

\subsubsection*{Was ist long?}\\
Der Datentyp \texttt{long} speichert sehr große ganze Zahlen, die über den \texttt{int}-Bereich hinausgehen.

\subsubsection*{Wertebereich:}
\begin{itemize}
	\item Von: -9.223.372.036.854.775.808
	\item Bis: 9.223.372.036.854.775.807
	\item Speichergröße: 64 Bit (8 Bytes)
\end{itemize}

\subsubsection*{Syntax:}
\begin{lstlisting}
long variablenName = wert;
long grosseZahl = wertL;  // L am Ende fuer long-Literal
\end{lstlisting}

\subsubsection*{Beispiele:}
\begin{lstlisting}
// Deklaration und Initialisierung
long weltbevoelkerung = 8000000000L;       // L am Ende!
long millisekunden = System.currentTimeMillis();

// Ohne L wird es als int interpretiert (Fehler bei großen Zahlen)
// long falsch = 8000000000;               // FEHLER!
long richtig = 8000000000L;                // Korrekt

// Berechnungen
long summe = 1000000000L + 2000000000L;
\end{lstlisting}

\begin{achtung}
	\begin{itemize}
		\item Bei großen Zahlen: \texttt{L} am Ende anhängen: \texttt{1000000000L}
		\item Ohne \texttt{L} wird Zahl als \texttt{int} interpretiert (kann zu Fehlern führen)
		\item Standard-Wert bei Arrays: \texttt{0L}
		\item Verwenden, wenn Zahlen größer als ca. 2 Milliarden sind
	\end{itemize}
\end{achtung}

% ============================================================
\newpage


\subsection*{float - Einfache Fließkommazahlen}
\addcontentsline{toc}{subsection}{float - Einfache Fließkommazahlen}

\subsubsection*{Was ist float?}\\
Der Datentyp \texttt{float} speichert Fließkommazahlen mit einfacher Genauigkeit.
Weniger genau als \texttt{double}.

\subsubsection*{Wertebereich:}
\begin{itemize}
	\item Ca. 6-7 Dezimalstellen Genauigkeit
	\item Von: $\pm$1.4 $\times$ 10$^{-45}$
	\item Bis: $\pm$3.4 $\times$ 10$^{38}$
	\item Speichergröße: 32 Bit (4 Bytes)
\end{itemize}

\subsubsection*{Syntax:}
\begin{lstlisting}
float variablenName = wertF;  // F am Ende fuer float-Literal
\end{lstlisting}

\subsubsection*{Beispiele:}
\begin{lstlisting}
// Deklaration und Initialisierung
float pi = 3.14F;                          // F am Ende!
float temperatur = 36.5F;
float gewicht = 75.8F;

// Ohne F wird es als double interpretiert
// float falsch = 3.14;                    // Warnung/Fehler
float richtig = 3.14F;                     // Korrekt

// Berechnungen
float summe = 10.5F + 20.3F;
\end{lstlisting}

\begin{achtung}
	\begin{itemize}
		\item \texttt{F} oder \texttt{f} am Ende anhängen: \texttt{3.14F}
		\item Weniger genau als \texttt{double}
		\item In der Praxis wird meist \texttt{double} bevorzugt
		\item Standard-Wert bei Arrays: \texttt{0.0F}
	\end{itemize}
\end{achtung}

% ============================================================
\newpage


\section*{String - Zeichenketten}
\addcontentsline{toc}{section}{String - Zeichenketten}

\subsection*{Was ist ein String?}

\subsubsection*{Beschreibung}\\
Ein \texttt{String} ist \textbf{kein primitiver Datentyp}, sondern eine Klasse (Referenztyp).
Er speichert Zeichenketten (Text).

\subsubsection*{Eigenschaften:}
\begin{itemize}
	\item Unveränderlich (immutable): Einmal erstellt, kann der Inhalt nicht geändert werden
	\item Wird mit \textbf{doppelten} Anführungszeichen erstellt: \texttt{"Text"}
	\item Kann beliebig lang sein
\end{itemize}

\subsubsection*{Syntax:}
\begin{lstlisting}
String variablenName = "Text";
\end{lstlisting}

\subsubsection*{Beispiele:}
\begin{lstlisting}
// Deklaration und Initialisierung
String name = "Max Mustermann";
String begruessung = "Hallo Welt!";
String leer = "";                          // Leerer String
String mitZahl = "Ich bin 25 Jahre alt";

// String-Konkatenation (Verkettung)
String vorname = "Max";
String nachname = "Mustermann";
String vollername = vorname + " " + nachname;  // "Max Mustermann"

// Mit Zahlen verketten
int alter = 25;
String text = "Ich bin " + alter + " Jahre alt";

// String-Methoden
String wort = "Hallo";
int laenge = wort.length();                // 5
char erstes = wort.charAt(0);              // 'H'
String gross = wort.toUpperCase();         // "HALLO"
String klein = wort.toLowerCase();         // "hallo"

// Strings vergleichen
String a = "Hallo";
String b = "Hallo";
boolean gleich = a.equals(b);              // true
boolean gleichIgnoreCase = a.equalsIgnoreCase("HALLO");  // true

// WICHTIG: Nicht mit == vergleichen!
String s1 = "Test";
String s2 = new String("Test");
boolean falsch = (s1 == s2);               // false (vergleicht Referenzen)
boolean richtig = s1.equals(s2);           // true (vergleicht Inhalt)

// Teilstrings
String text2 = "Hallo Welt";
String teil = text2.substring(0, 5);       // "Hallo"
String ab = text2.substring(6);            // "Welt"

// Suchen in Strings
String satz = "Java ist toll";
boolean enthaelt = satz.contains("ist");   // true
int position = satz.indexOf("ist");        // 5
boolean beginnt = satz.startsWith("Java"); // true
boolean endet = satz.endsWith("toll");     // true

// Ersetzen
String original = "Hallo Welt";
String neu = original.replace("Welt", "Java");  // "Hallo Java"

// Leerzeichen entfernen
String mitLeer = "  Text  ";
String ohneLeer = mitLeer.trim();          // "Text"

// String in Array umwandeln
String liste = "Apfel,Birne,Orange";
String[] fruechte = liste.split(",");      // ["Apfel", "Birne", "Orange"]

// Zahlen zu String
int zahl = 42;
String zahlString = String.valueOf(zahl);  // "42"
// ODER
String zahlString2 = "" + zahl;            // "42"

// String zu Zahl
String nummerText = "123";
int nummer = Integer.parseInt(nummerText); // 123
double komma = Double.parseDouble("3.14"); // 3.14
\end{lstlisting}

\begin{achtung}
	\begin{itemize}
		\item \textbf{Doppelte} Anführungszeichen: \texttt{"Text"}, nicht \texttt{'Text'}
		\item Strings mit \texttt{.equals()} vergleichen, \textbf{nicht} mit \texttt{==}
		\item \texttt{String} mit großem \textbf{S} schreiben
		\item Strings sind unveränderlich: Methoden geben neuen String zurück
		\item \texttt{null} ist möglich, aber \texttt{""} (leerer String) ist etwas anderes
		\item Standard-Wert bei Arrays: \texttt{null}
	\end{itemize}
\end{achtung}

% ============================================================
\newpage


\section*{Typkonvertierung}
\addcontentsline{toc}{section}{Typkonvertierung}

\subsection*{Implizite Konvertierung (Automatisch)}

\subsubsection*{Beschreibung}\\
Java konvertiert automatisch von \("\)kleinerem\("\) zu \("\)größerem\("\) Datentyp ohne Datenverlust.

\subsubsection*{Beispiele:}
\begin{lstlisting}
// int zu double (automatisch)
int ganzzahl = 5;
double kommazahl = ganzzahl;               // kommazahl = 5.0

// byte zu int (automatisch)
byte klein = 100;
int gross = klein;                         // gross = 100

// char zu int (automatisch)
char buchstabe = 'A';
int code = buchstabe;                      // code = 65 (Unicode-Wert)

// In Berechnungen
int a = 5;
double b = 2.5;
double ergebnis = a + b;                   // ergebnis = 7.5
\end{lstlisting}

\subsection*{Explizite Konvertierung (Type Cast)}

\subsubsection*{Beschreibung}\\
Von \("\)größerem\("\) zu \("\)kleinerem\("\) Datentyp muss explizit gecastet werden.
Kann zu Datenverlust führen!

\subsubsection*{Syntax:}
\begin{lstlisting}
zielTyp variable = (zielTyp) quellVariable;
\end{lstlisting}

\subsubsection*{Beispiele:}
\begin{lstlisting}
// double zu int (Nachkommastellen gehen verloren!)
double kommazahl = 3.7;
int ganzzahl = (int) kommazahl;            // ganzzahl = 3 (nicht 4!)

// long zu int
long gross = 1000L;
int klein = (int) gross;

// int zu byte (Vorsicht bei großen Zahlen!)
int zahl = 200;
byte kleinByte = (byte) zahl;              // Overflow!

// Division mit Cast
int a = 5;
int b = 2;
double ergebnis = (double) a / b;          // ergebnis = 2.5

// Runden statt Abschneiden
double wert = 3.7;
int gerundet = (int) Math.round(wert);     // gerundet = 4
\end{lstlisting}

\begin{achtung}
	\begin{itemize}
		\item Cast von double zu int: Nachkommastellen werden \textbf{abgeschnitten}, nicht gerundet!
		\item Bei Überschreitung: Datenverlust oder unerwartete Werte
		\item Für Rundung: \texttt{Math.round()} verwenden
		\item Implizite Konvertierung: klein → groß
		\item Explizite Konvertierung (Cast): groß → klein
	\end{itemize}
\end{achtung}

% ============================================================
\newpage


\section*{Übersicht: Primitive Datentypen}
\addcontentsline{toc}{section}{Übersicht: Primitive Datentypen}

\begin{table}[h]
	\centering
	\begin{tabular}{|l|l|l|l|}
		\hline
		\textbf{Typ} & \textbf{Größe} & \textbf{Wertebereich} & \textbf{Standard} \\
		\hline
		byte & 8 Bit & -128 bis 127 & 0 \\
		\hline
		short & 16 Bit & -32.768 bis 32.767 & 0 \\
		\hline
		int & 32 Bit & -2.147.483.648 bis 2.147.483.647 & 0 \\
		\hline
		long & 64 Bit & ca. -9,2 $\times$ 10$^{18}$ bis 9,2 $\times$ 10$^{18}$ & 0L \\
		\hline
		float & 32 Bit & ca. $\pm$3,4 $\times$ 10$^{38}$ & 0.0F \\
		\hline
		double & 64 Bit & ca. $\pm$1,7 $\times$ 10$^{308}$ & 0.0 \\
		\hline
		char & 16 Bit & 0 bis 65.535 (Unicode) & \('\)\textbackslash u0000' \\
		\hline
		boolean & 1 Bit & true oder false & false \\
		\hline
	\end{tabular}
	\caption{Übersicht der primitiven Datentypen in Java}
\end{table}


\section*{Arrays}
\addcontentsline{toc}{section}{Arrays}

\subsection*{Was ist ein Array?}

\subsubsection*{Beschreibung}\\
Ein Array ist eine Datenstruktur, die eine \textbf{feste Anzahl} von Elementen des \textbf{gleichen Datentyps} speichert.
Arrays sind Referenztypen (wie Strings), keine primitiven Datentypen.

\subsubsection*{Eigenschaften:}
\begin{itemize}
	\item Feste Größe: Nach der Erstellung kann die Größe nicht mehr geändert werden
	\item Index beginnt bei 0: Erstes Element hat Index 0, letztes Element hat Index \texttt{length - 1}
	\item Alle Elemente müssen den gleichen Typ haben
	\item Standard-Werte: Arrays werden automatisch mit Standard-Werten initialisiert (0, false, null, etc.)
\end{itemize}

\subsubsection*{Syntax:}
\begin{lstlisting}
// Deklaration und Initialisierung mit Größe
datentyp[] arrayName = new datentyp[groesse];

// Deklaration und Initialisierung mit Werten
datentyp[] arrayName = {wert1, wert2, wert3};
\end{lstlisting}

\subsubsection*{Beispiele:}
\begin{lstlisting}
// Array mit fester Größe erstellen
int[] zahlen = new int[5];        // 5 Plätze, alle 0
String[] namen = new String[3];   // 3 Plätze, alle null
boolean[] flags = new boolean[4]; // 4 Plätze, alle false

// Array mit Werten erstellen
int[] noten = {1, 2, 3, 2, 1};
String[] wochentage = {"Mo", "Di", "Mi", "Do", "Fr"};
double[] preise = {9.99, 14.50, 7.25};

// Mehrdimensionale Arrays
int[][] matrix = new int[3][3];   // 3x3 Matrix
int[][] tabelle = {{1, 2}, {3, 4}, {5, 6}};
\end{lstlisting}

\begin{achtung}
	\begin{itemize}
		\item Größe ist \textbf{fix} nach Erstellung
		\item Index beginnt bei \textbf{0}, nicht bei 1
		\item Alle Elemente müssen den \textbf{gleichen Typ} haben
		\item Standard-Wert bei Array-Variablen: \texttt{null}
	\end{itemize}
\end{achtung}

% ============================================================
\newpage


\subsection*{array.length - Länge eines Arrays}
\addcontentsline{toc}{subsection}{array.length - Länge eines Arrays}

\subsubsection*{Was ist array.length?}\\
\texttt{length} ist ein \textbf{Attribut} (keine Methode!), das die Anzahl der Elemente im Array zurückgibt.

\subsubsection*{Eigenschaften:}
\begin{itemize}
	\item Ist ein Attribut, keine Methode: \textbf{ohne} Klammern \texttt{()}
	\item Gibt die Größe zurück, mit der das Array erstellt wurde
	\item Funktioniert auch bei leeren Arrays (nur Standard-Werte)
\end{itemize}

\subsubsection*{Syntax:}
\begin{lstlisting}
int laenge = arrayName.length;  // OHNE Klammern!
\end{lstlisting}

\subsubsection*{Beispiele:}
\begin{lstlisting}
// Länge verschiedener Arrays
int[] zahlen = {10, 20, 30, 40, 50};
int anzahl = zahlen.length;               // anzahl = 5

String[] namen = new String[10];
int groesse = namen.length;               // groesse = 10 (auch wenn leer!)

// Leeres Array
int[] leer = new int[0];
int laenge = leer.length;                 // laenge = 0

// In Schleife verwenden
int[] werte = {1, 2, 3, 4};
for (int i = 0; i < werte.length; i++) {
    System.out.println(werte[i]);
}

// Letztes Element zugreifen
int[] liste = {100, 200, 300};
int letzter = liste[liste.length - 1];   // letzter = 300

// Bei mehrdimensionalen Arrays
int[][] matrix = new int[3][4];
int zeilen = matrix.length;               // zeilen = 3
int spalten = matrix[0].length;           // spalten = 4
\end{lstlisting}

\begin{achtung}
	\begin{itemize}
		\item \texttt{.length} \textbf{ohne Klammern} (bei Arrays)
		\item \texttt{.length()} \textbf{mit Klammern} (bei Strings)
		\item Gibt die Größe zurück, mit der das Array erstellt wurde
		\item Auch bei "leeren" Arrays (nur \texttt{null}-Werte) ist die Länge die Anzahl der Plätze
	\end{itemize}
\end{achtung}

% ============================================================
\newpage


\subsection*{array[index] - Element lesen}
\addcontentsline{toc}{subsection}{array[index] - Element lesen}

\subsubsection*{Was ist array[index]?}\\
Mit der eckigen Klammer \texttt{[]} greift man auf ein einzelnes Element im Array zu.
Der Index gibt die Position an.

\subsubsection*{Eigenschaften:}
\begin{itemize}
	\item Index beginnt bei 0: Erstes Element ist \texttt{array[0]}
	\item Letztes Element ist \texttt{array[array.length - 1]}
	\item Bei ungültigem Index: \texttt{ArrayIndexOutOfBoundsException}
\end{itemize}

\subsubsection*{Syntax:}
\begin{lstlisting}
datentyp element = arrayName[index];
\end{lstlisting}

\subsubsection*{Beispiele:}
\begin{lstlisting}
// Einzelne Elemente lesen
int[] zahlen = {10, 20, 30, 40, 50};
int erstes = zahlen[0];                   // erstes = 10
int zweites = zahlen[1];                  // zweites = 20
int letztes = zahlen[4];                  // letztes = 50

// Alle Elemente durchlaufen
String[] woerter = {"Hallo", "Welt", "Java"};
for (int i = 0; i < woerter.length; i++) {
    String wort = woerter[i];
    System.out.println(wort);
}

// Bestimmte Position prüfen
int[] noten = {1, 2, 3, 2, 1};
if (noten[2] == 3) {
    System.out.println("Dritte Note ist eine 3");
}

// Letztes Element
int[] liste = {5, 10, 15, 20};
int letzterWert = liste[liste.length - 1];  // letzterWert = 20

// Element in Berechnung verwenden
int[] preise = {100, 200, 150};
int summe = preise[0] + preise[1] + preise[2];  // summe = 450

// Mehrdimensionale Arrays
int[][] matrix = {{1, 2, 3}, {4, 5, 6}};
int wert = matrix[0][1];                  // wert = 2 (1. Zeile, 2. Spalte)
\end{lstlisting}

\begin{achtung}
	\begin{itemize}
		\item Index beginnt bei \textbf{0}, nicht bei 1
		\item Index muss zwischen 0 und \texttt{array.length - 1} liegen
		\item Bei ungültigem Index: \texttt{ArrayIndexOutOfBoundsException}
		\item Letztes Element: \texttt{array[array.length - 1]}
	\end{itemize}
\end{achtung}

% ============================================================
\newpage


\subsection*{array[index] = wert - Element schreiben}
\addcontentsline{toc}{subsection}{array[index] = wert - Element schreiben}

\subsubsection*{Was ist array[index] = wert?}\\
Mit der Zuweisung kann man den Wert eines Elements an einer bestimmten Position im Array ändern.

\subsubsection*{Eigenschaften:}
\begin{itemize}
	\item Überschreibt den vorherigen Wert an dieser Position
	\item Der neue Wert muss zum Datentyp des Arrays passen
	\item Index muss gültig sein (0 bis \texttt{length - 1})
\end{itemize}

\subsubsection*{Syntax:}
\begin{lstlisting}
arrayName[index] = neuerWert;
\end{lstlisting}

\subsubsection*{Beispiele:}
\begin{lstlisting}
// Array befüllen
int[] zahlen = new int[5];                // Array mit 5 Plätzen, alle 0
zahlen[0] = 10;                           // Erstes Element auf 10 setzen
zahlen[1] = 20;                           // Zweites Element auf 20 setzen
zahlen[4] = 50;                           // Letztes Element auf 50 setzen
// Ergebnis: {10, 20, 0, 0, 50}

// Werte ändern
String[] namen = {"Anna", "Bob", "Clara"};
namen[1] = "Max";                         // "Bob" wird zu "Max"
// Ergebnis: {"Anna", "Max", "Clara"}

// In Schleife befüllen
int[] quadrate = new int[5];
for (int i = 0; i < quadrate.length; i++) {
    quadrate[i] = i * i;
}
// Ergebnis: {0, 1, 4, 9, 16}

// Wert erhöhen
int[] punkte = {10, 20, 30};
punkte[0] = punkte[0] + 5;                // Erstes Element um 5 erhöhen
// ODER kürzer:
punkte[0] += 5;
// Ergebnis: {15, 20, 30}

// Array "leeren" (auf Standard-Wert setzen)
int[] werte = {1, 2, 3, 4, 5};
for (int i = 0; i < werte.length; i++) {
    werte[i] = 0;
}
// Ergebnis: {0, 0, 0, 0, 0}

// Zwei Elemente tauschen
int[] liste = {10, 20, 30};
int temp = liste[0];                      // Zwischenspeichern
liste[0] = liste[2];                      // 30 nach vorne
liste[2] = temp;                          // 10 nach hinten
// Ergebnis: {30, 20, 10}

// Array aus Benutzereingabe füllen
int[] zahlenListe = new int[3];
for (int i = 0; i < zahlenListe.length; i++) {
    // zahlenListe[i] = scanner.nextInt();
}

// Mehrdimensionale Arrays
int[][] matrix = new int[2][2];
matrix[0][0] = 1;
matrix[0][1] = 2;
matrix[1][0] = 3;
matrix[1][1] = 4;
// Ergebnis: {{1, 2}, {3, 4}}
\end{lstlisting}

\begin{achtung}
	\begin{itemize}
		\item Index muss zwischen 0 und \texttt{array.length - 1} liegen
		\item Bei ungültigem Index: \texttt{ArrayIndexOutOfBoundsException}
		\item Der Wert muss zum Datentyp des Arrays passen
		\item Überschreibt den vorherigen Wert an dieser Position komplett
	\end{itemize}
\end{achtung}

% ============================================================
\newpage


\section*{Häufige Fehler}
\addcontentsline{toc}{section}{Häufige Fehler}

\subsection*{String mit == vergleichen}

\begin{lstlisting}
// FALSCH
String a = "Hallo";
String b = "Hallo";
if (a == b) { }  // Vergleicht Referenzen, nicht Inhalt!

// RICHTIG
if (a.equals(b)) { }  // Vergleicht Inhalt
\end{lstlisting}

\subsection*{Integer-Division}

\begin{lstlisting}
// FALSCH - ergibt 2, nicht 2.5
int a = 5;
int b = 2;
double ergebnis = a / b;  // ergebnis = 2.0

// RICHTIG
double ergebnis = (double) a / b;  // ergebnis = 2.5
// ODER
double ergebnis = 5.0 / 2.0;  // ergebnis = 2.5
\end{lstlisting}

\subsection*{char vs. String verwechseln}

\begin{lstlisting}
// FALSCH
char buchstabe = "A";  // Fehler: doppelte Anfuehrungszeichen
String wort = 'Hallo';  // Fehler: einfache Anfuehrungszeichen

// RICHTIG
char buchstabe = 'A';   // einfache Anfuehrungszeichen
String wort = "Hallo";  // doppelte Anfuehrungszeichen
\end{lstlisting}

\subsection*{L und F vergessen}

\begin{lstlisting}
// FALSCH
long gross = 5000000000;  // Fehler: zu gross fuer int
float zahl = 3.14;        // Warnung: double wird zu float

// RICHTIG
long gross = 5000000000L;  // L am Ende
float zahl = 3.14F;        // F am Ende
\end{lstlisting}

\subsection*{Overflow nicht beachten}

\begin{lstlisting}
// Problem
byte b = 127;
b++;  // b = -128 (Overflow!)

int i = 2147483647;
i++;  // i = -2147483648 (Overflow!)
\end{lstlisting}

\subsection*{null vs. leerer String}

\begin{lstlisting}
String a = null;      // Keine Referenz (keine Methoden aufrufbar!)
String b = "";        // Leerer String (Methoden aufrufbar)

// a.length();        // FEHLER: NullPointerException
int laenge = b.length();  // Funktioniert: laenge = 0
\end{lstlisting}

\subsection*{String vs. Array - length}

\begin{lstlisting}
// FALSCH
String text = "Hallo";
int laenge = text.length;  // Fehler: Strings brauchen ()

int[] zahlen = {1, 2, 3};
int anzahl = zahlen.length();  // Fehler: Arrays haben KEINE ()

// RICHTIG
int laenge = text.length();    // String MIT Klammern
int anzahl = zahlen.length;    // Array OHNE Klammern
\end{lstlisting}

\subsection*{Array-Index-Grenzen}

\begin{lstlisting}
// FALSCH
int[] zahlen = {10, 20, 30};  // length = 3, Indizes: 0, 1, 2
int x = zahlen[3];            // Fehler! Index 3 existiert nicht
int y = zahlen[-1];           // Fehler! Negative Indizes nicht erlaubt

// RICHTIG
int letzter = zahlen[zahlen.length - 1];  // Index 2
int erster = zahlen[0];                   // Index 0
\end{lstlisting}

\subsection*{Array-Größe nach Erstellung ändern}

\begin{lstlisting}
// FALSCH
int[] zahlen = new int[5];
zahlen.length = 10;  // Fehler! length ist read-only

// RICHTIG - neues Array erstellen und Werte kopieren
int[] zahlen = new int[5];
int[] neuesArray = new int[10];
for (int i = 0; i < zahlen.length; i++) {
    neuesArray[i] = zahlen[i];
}
zahlen = neuesArray;  // Referenz auf neues Array
\end{lstlisting}
